\section{Related Work}
\label{sec:related}

\subsection{Statistical Watermarks for Language Models}
KGW \citep{kirchenbauer2023watermark} partitions the vocabulary with a keyed hash of the preceding context and biases green logits. \citet{kirchenbauer2024reliability} study its robustness and introduce additional seeding schemes, including the self-hashing (``selfhash'') scheme used in this work. Unigram watermarking fixes a single global partition \citep{zhao2024provable}. Distortion-free schemes preserve the next-token distribution in expectation over keys \citep{kuditipudi2024robust,mao2025unbiased}, and SynthID-Text uses tournament sampling at production scale \citep{dathathri2024scalable}. \citet{fernandez2023threebricks} sharpen the statistical tests used for detection, and \citet{jovanovic2024stealing} show that the keyed rules can be estimated and spoofed through queries. Quality is usually evaluated on aggregate text benchmarks \citep{tu2024waterbench,piet2023markmywords}. Some text-domain work reports variation in perplexity across keys \citep{mao2025unbiased}. \emph{Contrast:} these works treat the key as a secret to protect or as a source of noise to average over. We treat key choice itself as a treatment whose effect on downstream program correctness and security has to be measured.

\subsection{Watermarking Generated Code}
Code has low-entropy stretches where any bias can break syntax. SWEET \citep{lee2024sweet} therefore watermarks only high-entropy tokens and reports that plain KGW reduces functional pass rates on code benchmarks. \citet{suresh2024robust} show that semantics-preserving transformations such as renaming remove code watermarks. \emph{Contrast:} these studies measure the \emph{average} cost of watermarking on functional correctness, or the watermark's robustness, for a fixed key. They do not measure security outcomes, and they do not test whether valid keys differ from one another. Our protocol estimates both the average effect and between-key heterogeneity, using a combined secure-and-functional endpoint.

\subsection{Security Evaluation of LLM-Generated Code}
Early studies found that roughly 40\% of Copilot completions in CWE-targeted scenarios were vulnerable \citep{pearce2022asleep}. SecurityEval \citep{siddiq2022securityeval} and CyberSecEval \citep{bhatt2023cyberseceval} provide CWE-mapped prompts with static detection, and SVEN steers models toward secure code \citep{he2023sven}. CWEval \citep{peng2025cweval} argues that static analyzers give unstable security judgments and that functionality and security should be evaluated on the same problem with outcome-driven tests. We adopt this position: our primary endpoint counts a program as a success only if it is both functional and secure. \emph{Contrast:} these benchmarks measure model-level security without watermarking. Our contribution is a paired design that isolates the watermark and its key as the only varied factor.
